\clearpage
\section{SneakySnake prefiltering}
\label{sec:prefilter-role}
\label{sec:prefilter-sneakysnake}

A prefilter rejects candidates with too many edits before base-level alignment.
BioseqDB uses SneakySnake~\cite{sneakysnake_paper,sneakysnake_readme}, which
models pairwise character comparisons as a \emph{chip maze} and searches
for a signal net through it.

For equal-length sequences \(Q,R\in\Sigma^m\) and edit limit \(E\), use the
paper's row order and columns \(1\leq j\leq m\):
\[
 \delta_i=\begin{cases}
  -i,&1\leq i\leq E,\\
  0,&i=E+1,\\
  i-E-1,&E+1<i\leq2E+1.
 \end{cases}
 \qquad
 Z_{i,j}=\begin{cases}
  0,&1\leq j+\delta_i\leq m\ \land\ Q_{j+\delta_i}=R_j,\\
  1,&\text{otherwise}.
 \end{cases}
\]
Each row is a horizontal track with zeros for matches and obstacles for
mismatches or out-of-range comparisons. From each checkpoint the net follows
the longest zero run across all rows, crosses the next obstacle if present,
and repeats. Row changes are free, so the obstacle count is a lower bound
on edit distance. Reject only if it exceeds \(E\).

\begin{figure}[H]
\centering
\begin{tikzpicture}[x=1cm,y=-1cm,>=Latex,font=\small,
  wire/.style={blue!65!black,line width=.45pt},
  net/.style={red!75!black,line width=1.1pt},
  obstacle/.style={fill=black!80,draw=white,line width=.3pt}]
  \node[anchor=west,font=\bfseries\small] at (-.15,-.2) {(a) Construct column 6};
  \node[anchor=west,font=\bfseries\small] at (6.6,-.2) {(b) Chip maze and signal net};

  % One reference symbol fans out to the seven query symbols Q_3,...,Q_9.
  \foreach \base [count=\j] in {G,G,T,G,C,A,G,A,G,C,T,C} {
    \node[font=\scriptsize,text=gray] at ({.75+.37*\j},.27) {\j};
    \node[font=\ttfamily\small] (ref\j) at ({.75+.37*\j},.68) {\base};
  }
  \node[anchor=east] at (.8,.68) {$R$};
  \draw[blue!65!black,thick,rounded corners=1pt]
    ($(ref6.center)+(-.17,-.2)$) rectangle ($(ref6.center)+(.17,.2)$);
  \foreach \base [count=\j] in {G,G,T,G,A,G,A,G,T,T,G,T} {
    \ifnum\j>2\ifnum\j<10
      \fill[blue!8] ({.75+.37*\j-.17},1.65) rectangle ({.75+.37*\j+.17},2.05);
    \fi\fi
    \node[font=\ttfamily\small] (query\j) at ({.75+.37*\j},1.85) {\base};
  }
  \node[anchor=east] at (.8,1.85) {$Q$};
  \foreach \j/\level in {3/1.06,4/1.18,5/1.30,6/1.42,7/1.30,8/1.18,9/1.06} {
    \draw[wire] (ref6.south) -- (2.97,\level) -- ({.75+.37*\j},\level) -- (query\j.north);
  }
  \node[font=\scriptsize] at (.1,2.28) {$i$};
  \node[font=\scriptsize] at (.65,2.28) {$Z_{i,6}$};
  % The first three results are reversed, exactly as in the paper's Eq. (1).
  \foreach \i/\j/\bit in {1/5/0,2/4/1,3/3/1,4/6/1,5/7/0,6/8/1,7/9/1} {
    \pgfmathsetmacro{\rowy}{2.72+.47*(\i-1)}
    \draw[wire,preaction={draw=white,line width=1.6pt}]
      (query\j.south) -- ({.75+.37*\j},\rowy) -- (.86,\rowy);
    \node[font=\scriptsize] at (.1,\rowy) {\i};
    \ifnum\bit=1
      \node[fill=black!80,text=white,minimum size=.31cm,inner sep=0pt,font=\scriptsize]
        at (.65,\rowy) {1};
    \else
      \node[draw=teal!65!black,fill=teal!8,minimum size=.31cm,inner sep=0pt,font=\scriptsize]
        at (.65,\rowy) {0};
    \fi
  }
  \node[font=\scriptsize,text=blue!65!black] at (3.25,5.95)
    {$R_6=\mathtt A,\quad Z_{\cdot,6}=(0,1,1,1,0,1,1)^{\mathsf T}$};

  % Full 7 x 12 maze; bit rows use delta = -1,-2,-3,0,+1,+2,+3.
  \begin{scope}[shift={(7.05,2.72)},x=.48cm,y=-.47cm]
    \fill[blue!10] (4.66,-.45) rectangle (5.34,6.45);
    \foreach \j in {0,...,11} {
      \draw[gray!40,densely dashed,thin] (\j,-.35) -- (\j,6.35);
      \node[font=\scriptsize] at (\j,-.8) {\pgfmathparse{int(\j+1)}\pgfmathresult};
    }
    \foreach \i in {0,...,6} {
      \draw[gray!45,line width=1.5pt] (-.4,\i) -- (11.4,\i);
      \node[font=\scriptsize,anchor=east] at (-.55,\i) {\pgfmathparse{int(\i+1)}\pgfmathresult};
    }
    \node[font=\scriptsize] at (-1.1,-.8) {$i\backslash j$};
    \foreach \row/\bits in {
      0/{1,0,1,1,1,0,0,0,0,1,0,1},
      1/{1,1,1,0,1,1,1,1,1,1,0,1},
      2/{1,1,1,0,1,1,0,0,0,1,1,1},
      3/{0,0,0,0,1,1,1,1,1,1,1,1},
      4/{0,1,1,1,1,0,0,1,1,1,0,1},
      5/{1,0,1,0,1,1,1,1,0,1,1,1},
      6/{0,1,1,1,1,1,1,1,1,1,1,1}} {
      \foreach \bit [count=\col from 0] in \bits {
        \ifnum\bit=1
          \path[obstacle] (\col-.17,\row-.17) rectangle (\col+.17,\row+.17);
        \fi
      }
    }
    % Escape segments cross obstacles at (row,column) = (4,5), (1,10), (1,12).
    \draw[net,->] (-.6,3) -- (4.42,3);
    \draw[net] (4.42,3) -- (4.42,0);
    \draw[net,->] (4.42,0) -- (9.42,0);
    \draw[net,->] (9.42,0) -- (11.55,0);
    \foreach \col/\row in {4/3,9/0,11/0} {
      \draw[orange!85!black,line width=.9pt] (\col,\row) circle[radius=.25];
    }
    \foreach \col/\label in {-.55/1,4.48/6,9.48/11} {
      \draw[teal!65!black,densely dashed] (\col,-1.55) -- (\col,6.4);
      \node[font=\scriptsize,text=teal!65!black,anchor=south] at (\col,-1.65) {$c=\label$};
    }
  \end{scope}
  \draw[net,->] (6.8,6.03) -- (7.5,6.03);
  \node[anchor=west,font=\scriptsize] at (7.65,6.03) {signal net};
  \path[obstacle] (9.45,5.94) rectangle (9.63,6.12);
  \node[anchor=west,font=\scriptsize] at (9.78,6.03) {obstacle};
  \draw[teal!65!black,densely dashed] (11.55,5.88) -- (11.55,6.18);
  \node[anchor=west,font=\scriptsize] at (11.7,6.03) {checkpoint};
\end{tikzpicture}
\caption{Reconstruction of Fig.~2(a,b) in
\citeauthor{sneakysnake_paper} (\citeyear{sneakysnake_paper}) for
\(R=\mathtt{GGTGCAGAGCTC}\), \(Q=\mathtt{GGTGAGAGTTGT}\), and \(E=3\).
(a) Comparing \(R_6\) with \(Q_3,\ldots,Q_9\) constructs column~6
(blue in (b)); the first three results follow the paper's reversed order.
(b) The net crosses three circled obstacles, so the pair passes the filter.
Only runs up to the first obstacle after each checkpoint need evaluation.}
\label{fig:prefilter-chip-maze}
\end{figure}

BioseqDB applies this test to mandatory intervals between anchors;
optional outer extensions are left to the aligner. For unequal lengths,
the implementation uses only \(\bigl||Q|-|T|\bigr|\).
Survivors are aligned and checked against the original edit predicate:
passing the filter alone does not establish that it holds.

% Existing filter illustrations retained as source material.
% \begin{figure}[H]
% \centering
% \begingroup
% \small
% \setlength{\tabcolsep}{4pt}
% \renewcommand{\arraystretch}{1.2}
% \newcommand{\mazeMatch}{\cellcolor{white}\(\cdot\)}
% \newcommand{\mazeObstacle}{\cellcolor{black!70}\textcolor{white}{\(\times\)}}
% \newcommand{\mazeRoute}{\cellcolor{orange!35}\textcolor{orange!80!black}{\(\rightarrow\)}}
% \newcommand{\mazeCross}{\cellcolor{orange!70}\(\times\)}
% \begin{tabular}{c|*{12}{c}}
%  & 0 & 1 & 2 & 3 & 4 & 5 & 6 & 7 & 8 & 9 & 10 & 11\\
% \hline
% \(d=+2\) & \mazeMatch & \mazeMatch & \mazeMatch & \mazeObstacle &
% \mazeMatch & \mazeMatch & \mazeMatch & \mazeObstacle & \mazeMatch &
% \mazeMatch & \mazeObstacle & \mazeObstacle\\
% \(d=+1\) & \mazeMatch & \mazeMatch & \mazeMatch & \mazeObstacle &
% \mazeMatch & \mazeMatch & \mazeMatch & \mazeObstacle & \mazeMatch &
% \mazeMatch & \mazeMatch & \mazeObstacle\\
% \(d=0\) & \mazeRoute & \mazeRoute & \mazeRoute & \mazeCross &
% \mazeRoute & \mazeRoute & \mazeRoute & \mazeCross & \mazeRoute &
% \mazeRoute & \mazeRoute & \mazeCross\\
% \(d=-1\) & \mazeObstacle & \mazeMatch & \mazeMatch & \mazeObstacle &
% \mazeMatch & \mazeMatch & \mazeMatch & \mazeObstacle & \mazeMatch &
% \mazeMatch & \mazeMatch & \mazeObstacle\\
% \(d=-2\) & \mazeObstacle & \mazeObstacle & \mazeMatch & \mazeObstacle &
% \mazeMatch & \mazeMatch & \mazeMatch & \mazeObstacle & \mazeMatch &
% \mazeMatch & \mazeMatch & \mazeObstacle\\
% \end{tabular}

% \vspace{0.4em}
% \begin{tabular}{clclcl}
% \cellcolor{white}\(\cdot\) & match &
% \cellcolor{black!70}\textcolor{white}{\(\times\)} & obstacle &
% \cellcolor{orange!35}\(\rightarrow\) & selected escape segment
% \end{tabular}
% \endgroup
% \caption{A rejecting SneakySnake chip-maze walk for
% \(Q=\mathtt{AAAAAAAAAAAA}\), \(T=\mathtt{AAACAAACAAAC}\), and \(E=2\).
% Each row represents a query-to-target displacement. Amber cells show the
% successive escape segments and three crossed obstacles on the unshifted row.}
% \label{fig:prefilter-chip-maze}
% \end{figure}

% \begin{figure}[H]
% \centering
% \small
% \renewcommand{\arraystretch}{1.15}
% \begin{tabular}{c}
% \fbox{\parbox{0.74\textwidth}{\centering
% Construct the unshifted mismatch mask \(H_0\)}}\\[0.35em]
% \(\operatorname{popcount}(H_0)\leq E\): retain \qquad
% \(E=0\): reject by exact Hamming comparison\\[0.35em]
% \(\Downarrow\) otherwise\\[0.35em]
% \fbox{\parbox{0.74\textwidth}{\centering
% Construct the \(2E\) shifted masks, amend each mask independently,\\
% then set the boundary-adjacent compensation bits}}\\[0.35em]
% \(\Downarrow\)\\[0.35em]
% \fbox{\parbox{0.74\textwidth}{\centering
% Intersect the masks and compute the capped four-bit group estimate}}\\[0.35em]
% \(\Downarrow\) retain when the estimate is at most \(E\)\\[0.35em]
% \fbox{\parbox{0.74\textwidth}{\centering
% For a proposed rejection, evaluate exact thresholded edit distance}}\\[0.35em]
% \(D(Q,T)\leq E\): retain \qquad \(D(Q,T)>E\): reject
% \end{tabular}
% \caption{GateKeeper decision flow in BioSeqDB. The unshifted mask provides an
% exact rejection certificate when \(E=0\). For positive thresholds, shifted
% masks can retain a candidate or propose rejection, and the exact threshold
% guard confirms the proposal.}
% \label{fig:prefilter-gatekeeper}
% \end{figure}
